% 从旧数学附录迁出的树路径说明；不新增算法。
\section{树路径交：投影与单点交}\label{knowledge-path-intersection}
\index{树上路径交}
核心见第~\pageref{compact-path_intersection}~页，\href{https://www.luogu.com.cn/problem/P3398}{P3398} 完整用法见第~\pageref{usage-example-169}~页。
将 $x$ 投影到树路径 $a$--$b$，投影是 $\operatorname{lca}(a,b)$、$\operatorname{lca}(a,x)$、$\operatorname{lca}(b,x)$ 中深度最大者。
三条两点间路径在同一个分叉点相遇；按同一根计算 LCA 即得这个分叉点，与顶点编号无关。
设 $c,d$ 的投影是 $p,q$。若 $p\ne q$，路径 $c$--$d$ 必经过 $p$--$q$，后者正是两路径的交集。
若 $p=q$，两端可能在同一条旁支，不能直接认为有交；须检查 $d(c,p)+d(p,d)=d(c,d)$。
通过时为单点交，否则为空。这里距离和 depth 都按边数，不能换成带权根距离；零权边会破坏上述距离等式的判别。

接口返回交集端点和顶点数；空交是 $\{0,0,0\}$，公共边数为 $\max(0,\mathrm{vertices}-1)$。
P3398 仅输出是否相交，不验证具体端点。
本地接口测试穷举至五点的全部标号树、每个根及四端点，独立构造两条路径的顶点集合求交，核验数量与返回端点间的整条路径。
完整驱动另验 146 棵小树、100 棵随机换号树，以及十万点链和十万次查询；两种编译模式均通过。
EulerLCA 预处理 $O(n\log n)$、每次路径交 $O(1)$；若换用 HLD 回调则为 $O(\log n)$。
完整驱动以 C++20 O2 在\href{https://www.luogu.com.cn/record/299972443}{P3398 记录 299972443}通过全部十个测试点，
最慢单点 134 ms、峰值 31.89 MB；该 AC 仅对应布尔相交判定，不涵盖端点、精确数量或带权距离。

\clearpage
\section{非负权树的点集直径摘要}\label{knowledge-set-diameter}
代码见第~\pageref{compact-TreeDiameter}~页。
一个非空点集只需保存任意一对直径端点 $a,b$，即可回答任意树上点 $x$ 到该点集的最远距离：
$\max(d(x,a),d(x,b))$。
边权非负，包括零权；该结论不能推广到一般图或负权树。

证明可将点 $u$ 投影到路径 $a$--$b$，记投影距 $a$ 为 $t$，离路径的距离为 $h$，
直径长为 $D$。由于 $d(a,u),d(b,u)\le D$，有 $h\le\min(t,D-t)$。
再将查询点投影到同一路径，投影距 $a$ 为 $s$，离路径的距离为 $g$。
经两个投影点连接得到
$d(x,u)\le g+h+|s-t|\le g+\max(s,D-s)$，
右侧恰为 $\max(d(x,a),d(x,b))$。
因此插入新点只需比较到两个旧端点的距离；合并点集时，只插入另一个摘要的两个端点即可。
并列直径可任选，不应假设端点编号固定或合并后的端点与合并顺序无关。

Euler LCA 现在同时保存根距离，但其原有 distance 接口仍返回边数。
带权距离请显式调用 weighted\_distance / Weighted\_Distance；
采用 $(D_u-D_z)+(D_v-D_z)$ 避免不必要的两倍根距离中间值。
所有路径权和须能用 long long 表示。
独立验证穷举小点集及合并，逐点检查最远距离与返回顶点；
带权路径用独立树遍历参照，并覆盖全部根、零边权与 \texttt{LLONG\_MAX} 单边。
原有 Euler LCA 全小树回归，以及 50 万点、边权 $10^9$ 的递归链也通过普通与 ASan/UBSan 检查。

\paragraph{离线添叶仍是固定树度量。}
\href{https://codeforces.com/problemset/problem/379/F}{CF379F New Year Tree} 的完整用法见第~\pageref{usage-example-170}~页。
先把全部添叶操作建成最终树，再用 HLD 求距离；只有旧树外接新叶，不会改变两个旧点之间的唯一简单路径。
因此可以按出现顺序向摘要插点，尚未出现的点只帮助预处理距离，不提前进入点集。
初始四点也必须插入；每次两个新点都插完，再输出 length。
$q\le500000$ 时最终有 $4+2q$ 个点，采用 HLD 的 $O(q)$ 空间、$O(q\log q)$ 总时间。
递归 DFS 保留；本地放大测试栈，最大实例包含五十万层链。
250 组随机合法添叶过程逐步重建当前树，以独立两次 BFS 求直径，另验满规模链；普通与 ASan/UBSan 均通过。
这份用法仅覆盖 insert 和 length；任意点集合并、空集和最远点见接口独立测试。
2026-09-29 洛谷提交入口返回 Codeforces RemoteJudge 暂不可用，故不登记此题线上 AC。

\href{https://qoj.ac/problem/8235}{杭州 2023 F} 中，点权互异。
令 $g$ 为全树点权的 mex，只需按点权 $0,1,\ldots,g-1$ 建立前缀摘要。
某个前缀全部位于以 $x$ 为中心、半径 $k$ 的范围内，当且仅当该摘要的最远距离不超过 $k$。
二分最长满足条件的前缀长度即为答案；$g=0$ 时直接为零。
题面树边带权，不能用边数距离。
完整驱动在 \texttt{verify/qoj/8235.*.cpp}，
本地使用实际权距与 mex 集合参照，$n=q=500000$ 的链和星形树保持递归 DFS。
两套驱动各 50 个正式测试和 1 个额外测试全部在线 AC：
\href{https://qoj.ac/submission/2939493}{动态版} 666 ms、198076 KB，
\href{https://qoj.ac/submission/2939551}{传统版} 588 ms、198196 KB。
OJ 验证题面正权边、前缀插入和最远距离判定；零权边、完整 64 位距离与任意摘要合并仍由独立本地验证支持。

\clearpage
